% !TeX root = ../papers/nc-p3.tex
\section{両対称反転窓の二つの経路長による必要十分判定}
\label{sec:bisymmetric-finite-path-test}
$R\subset V\otimes W$を隣接階数6の両対称Gramテンソルの関係六平面とし、
反転窓の二次代数を$C=C(R,R^t,R,R^t)$とする。以下の判定は
反対階数、中心、反対鉛筆の曲線型の事前判定を必要としない。
$P_d$を二変数$d$次形式の空間と書く。
\begin{theorem}\label{thm:bisymmetric-two-length-test}
次は同値である。
\begin{enumerate}
\item $C$は標準14641型Koszul左右AS正則である。
\item いずれの始点でも長さ3の経路次元が20以下、長さ21の経路次元が2024以下である。
\item $R$は非零の純テンソルを含まず、どちらの非零片側評価も階数3以上であり、
次に定める二変数形式への商がどちらの向きにも存在しない。
\end{enumerate}
ここで二変数形式への商とは、矢空間の全射$V\to P_2^*$, $W\to P_3^*$で
$R$の像が乗法$P_2\otimes P_3\to P_5$の双対像に入るものをいう。
最初の二つの線型代数条件（純関係なし、評価階数3以上）を直接確認した場合は、
経路長17で両始点とも1140以下であることだけで全次数の正則性を判定できる。
特に両対称の隣接階数6の空間で標準正則窓はZariski開集合をなす。
\end{theorem}
\begin{proof}
まず第三条件の十分性を示す。反対核に可約な元があれば
補題\ref{lem:bisymmetric-reducible-quick-test}による。以下、どちらの核にも
そのような元がないとする。定理\ref{thm:all-bisymmetric-as-closure}の証明で、
楕円鉛筆閉包に入るために使う条件は純関係の不在と両軸での評価階数だけであった。
従ってこの窓を、一般点で相異なる二束を持つ滑らかな楕円鉛筆から特殊化できる。
二側の定義二次式の二平面をGrassmann多様体で極限に取る。それぞれ実際の
反対核に含まれるので、共通の一次因子も階数2以下の元もない。
従って両方とも四次完全交叉であり、非被約な鉛筆には可約な二次式があるという
定理\ref{thm:nonreduced-quartic-pencils}の議論から被約である。
反対階数が特殊点で一定であると仮定する必要はない。

被約極限について、既出の特殊化証明がどの障害で場合を除くかを確認する。
滑らかな側があれば、評価階数により切断格子が揃い、純関係がないので基点自由な
鉛筆が延びる。整な節点・尖点の場合にも定理\ref{thm:integral-apolar-alternatives}
の格子のずれは実際の評価階数2以下を与える。格子が揃った後、可逆束なら
基点自由な整曲線の表示になる。非可逆なら導手で消える三切断と他側の全四切断の
積が$P_2\otimes P_3\to P_5$を経由し、禁止された商を与える。

両側が節点多角形の場合、定理\ref{thm:two-polygon-specialization}の証明で
多重次数の異なる172型は、極限の積の階数の損失を含めても実際の評価核が
二次元以上になる。残る二型には命題\ref{prop:semipositive-pencil-balance}の
二変数形式の商がある。多角形と楕円多重点を組み合わせる場合は
命題\ref{prop:polygon-additive-exclusion}の証明が、実際の評価核に二次元を与える。
最後に両側が多重点の場合、定理\ref{thm:all-reduced-apolar-specializations}の
証明では、中心対のない場合は積の階数が8を超えるか実際の評価階数が2以下となる。
中心対がある場合も、積の階数7の不一致および次数の大きい末端は同じ評価障害で落ちる。
残る共通多重点模型では積が八次元を張り、命題\ref{prop:semipositive-pencil-balance}
によって、次数不一致は評価障害か二変数形式の商のいずれかを与える。
これらは上記各証明で正則性の必要条件として使われた具体的な排除であり、
ここでは第三条件が直接これらを禁じている。残った等次数の模型は積の二次元鉛筆を
持ち、基点は純関係を与えるため存在しない。各被約曲線の全次数定理が適用できる。

第二条件から第三条件を導く。非零評価の核が二次元以上なら、一文字とその核の
二文字へ交互に射影して全関係を消せる。二文字始点の長さ21では
$2^{11}=2048>2024$となる。従って両評価は階数3以上である。
特に各一文字での最上テンソルの縮約は単射で、三次交叉に四次元の像を持つ。
長さ3の次元は$64-2\cdot24+\dim W_{i,3}$なので、この像は交叉に等しい。
もし$a\otimes b\in R$が純関係なら、転置により$a\otimes b\otimes a$が交叉に入る。
従ってこれはある非零文字による最上テンソルの縮約である。
その先頭・末尾の一次元支持は、その文字で隣接関係を評価した像に等しい
（定理\ref{thm:no-consecutive-pure-relations}の縮約の議論）。
評価階数が1となり矛盾する。

二変数形式の商があれば反転先でも転置により同じ商が続く。
補題\ref{lem:binary-diagonal-paths}の経路商の$P_3^*$始点で長さ$2m+1$の次元は
$9\cdot2^{m-1}$である。長さ21では4608となり第二条件に反する。
長さ17でも$1152>1140$なので追加の判定も従う。
第一条件なら標準Hilbert級数から第二条件を得る。最後に経路関係の挿入行列の
階数下界は、六次元関係束の各基底チャートで開条件である。
\end{proof}

二変数形式の商は曲線の同型を探索せずに検査できる。三行四列の独立行列$U$を取り、
関係の各係数行列$A$について$UAY$がHankel行列となる可逆な四行四列行列$Y$を探す。
六つの反対角線が一定となることは、相隣る二行のずれから各関係につき六つの
線型方程式を与える。固定した$U$に対し合計36行16列の核を取り、その線型系で
行列式が恒等零かどうかで可逆解を判定できる。$U$については三次小行列が一つ
非零となる四つのチャートで解の存在を調べればよく、左右を交換した検査も行う。
\path{computations/verify_bisymmetric_finite_test.py}は節点・尖点の二標本で
座標を取り替え、指定した三平面上のこの核から可逆な解をスカラーまで復元する。
